In modern digital systems, shift registers support data alignment, serialization, and buffering, so verification must remain correct and scalable as features such as shift direction, parallel load, and reset behavior are added. However, a simple directed SystemVerilog testbench (even with assertions) can miss corner cases and does not scale well when the design grows or when similar modules (ex: FIFOs) require verification. This project addresses that gap by building a reusable, constrained-random verification environment for a shift register using the Universal Verification Methodology (UVM) rather than a simple SystemVerilog testbench.
The key idea is to take advantage of UVM’s modular architecture. The Interface connects the DUT to the testbench, the driver drives signals into the DUT based on the sequencer’s transactions, the monitor observes signals and converts activities into transactions, and the scoreboard validates the correct output based on the predictor inside the UVM environment. This structure enables systematic verification of both parallel-to-serial and serial-to-parallel behaviors across many randomized scenarios, while keeping the environment reusable to future designs. A core challenge is defining transactions and a cycle-accurate predictor that models the shift behavior correctly, then ensuring the monitor samples signals at the right times to avoid timing issues during scoreboard comparisons.
The measurable outcome is coverage, which is collected by the covergroups/coverpoints. The goal is to achieve 100% functional coverage, demonstrating that all planned scenarios in the test plan are exercised and checked. Limitations include the effort required to design the testbench. A simple module such as a shift register, still needs the extensive UVM design in order to get a basic testbench up and running, while a SystemVerilog testbench is easier, and faster to set up. Next, this UVM framework will be adapted to verify a FIFO by reusing the same architectural components (sequences, monitors, and scoreboards), demonstrating how verification reuse reduces development time and increases confidence in larger digital subsystems.

\lipsum[1]